% !TEX program = xelatex
% 中国研究生数学建模竞赛 F 题论文初稿；数据与结果来自题目附件。
\documentclass[fontset=windows,12pt,a4paper]{ctexart}
\usepackage[a4paper,top=27mm,bottom=25mm,left=23mm,right=23mm]{geometry}
\usepackage{amsmath,amssymb,booktabs,array,graphicx,float,hyperref}
\usepackage{caption}
\hypersetup{colorlinks=false,pdfborder={0 0 0}}
\captionsetup{font=small}
\setlength{\parindent}{2em}
\linespread{1.35}
\graphicspath{{../outputs/}}
\newcommand{\fig}[3]{\begin{figure}[htbp]\centering\includegraphics[width=#2\linewidth]{#1}\caption{#3}\end{figure}}
\title{算力约束下大语言模型能力提升的资源配置建模\\\large 基于题目附件的质量、配比、标度与开放模型证据}
\author{}
\date{}
\begin{document}
\maketitle
\begin{abstract}
针对题目提出的数据质量评价、训练标度、资源配置与开放模型能力前沿四个问题，本文按附件证据类型建立可复算的分层模型，并明确跨附件不能识别的关系。问题一对 A1--A3 共 272505 条记录的 22 项指标构造四组稳健质量分，分析指标冲突；对 17 域训练配比采用二次 ILR 岭回归，独立 1M 测试的平均 Loss 均方根误差为 0.2319、Spearman 相关为 0.6659。推荐配比取 52 个实测优良配方的质心，不称其为全局最优。问题二基于 B1 的 1176 个检查点拟合参数量--数据量标度律，再用 B6 半合成数据分组验证质量修正结构；两项均受质量因子作用的候选模型取得 0.0595 的分组交叉验证误差。第一问质量分 $Q_A$ 与 B6 的 $Q_B$ 缺少成对标定，故不将二者直接代入同一数值式。问题三固定第一问配比候选，在观测支持域内求解不同预算的 $(N,D,Q_B)$；$10^{24}$ FLOPs 情景有 97.58\% 预算未用，原因是支持域上界而非实际高预算配置最优。问题四在严格开放筛选下分别分析预训练与后训练模型；现有可比数据不能支持可靠的 12/24 个月能力分数预测，因此只给出题目附件支持的条件算力路径与能力识别边界。

\noindent\textbf{关键词：}数据质量；组成数据；标度律；算力配置；部分识别；开放模型
\end{abstract}
\tableofcontents
\newpage
\section{问题重述、数据分层与符号}
\subsection{任务与递进关系}
题目与随附数据说明\cite{problem}要求依次回答四项任务：从多指标评价预训练文本质量并分析冲突、建立领域配比与训练损失的关系；结合模型参数量 $N$、训练数据量 $D$、质量 $Q$ 与配比 $p$ 讨论标度律和边际效用；在给定算力预算下配置资源；最后分析开放模型能力变化和算力增速放缓的情景。本文以题目附件为主证据。四问按“质量和配比候选$\rightarrow$标度响应$\rightarrow$预算配置$\rightarrow$能力边界”的顺序连接，但只对有共同观测协议的量作数值传递。

\subsection{数据角色及不可混用的标尺}
A1--A3 是质量指标记录，A4--A15 提供配比及不同规模的 Loss，A16 为领域映射参考；B1 为训练检查点，B3 是插值轨迹，B6--B8 涉及半合成质量实验；C1--C8 包含开放模型排行榜、权重与许可、算力及任务信息。B3 插值记录不用于独立验证；A12--A15 的大尺度 Loss 估计不当成真实训练观测。$Q_A$ 表示问题一的相对质量评分，$Q_B$ 表示 B6 的原生质量水平。两者没有相同样本的成对校准数据，不能直接互换。不同来源的 Loss 也不强行拼接为同一训练集。

本文令 $N_B=N/10^9$、$D_B=D/10^9$，即参数量和 token 数均以十亿计；$C$ 以 FLOPs 计。$p=(p_1,\ldots,p_{17})$ 为非负且和为 1 的领域配比。$L$ 表示各附件自身协议下的 Loss；排行榜六任务均分以分数点计，与 $L$ 不同量纲。

\begin{table}[htbp]\centering\caption{四问主要证据及可回答范围}
\small\begin{tabular}{p{1.3cm}p{3.1cm}p{7.3cm}}\toprule
问题 & 主证据 & 可回答范围\\\midrule
一 & A1--A7、A16 & 数据内相对质量、1M 配比预测与候选配方\\
二 & B1、B6；B4/B5/B7 验证 & 标度关系、半合成质量效应和条件边际量\\
三 & 第一问候选配比、第二问响应、题设成本 & 指定支持域与成本函数下的条件配置\\
四 & C1/C2/C4/C6/C8 & 开放样本的描述性趋势、算力情景与预测可识别性\\\bottomrule
\end{tabular}\end{table}

\subsection{建模原则}
先核对文件编号、行数与字段，再划分训练、验证、插值和估算数据；同一模型家族或训练轨迹整体留出，避免相邻记录泄漏。所有模型结论附上支持域、误差口径和证据类型。题目允许对配比采用前问结果作为固定方案，本文在第三问采用这一选择，并明确其尚未经跨问联合实验检验。

\section{问题一：质量评价与领域配比}
\subsection{质量指标统一与冲突分析}
A1--A3 合计 272505 条记录。对 A1 的 22 项指标按题给语义统一正向与反向，在 A1 的 $1\%$--$99\%$ 分位确定稳健标尺，再原样用于 A2/A3。将指标分为知识价值、语言表达、结构丰富和洁净规范四组。每组先求稳健组分，再用 Huber 位置函数合成相对质量分 $Q_A$；其作用是比较当前附件中的样本，不能解释成单篇文本的真实训练收益。冲突度定义为四组评分六个两两绝对差的平均值；最大的差值对仅用于标注主导冲突方向，以 A1 的第 90 百分位 $0.2714$ 为阈值，随后保持阈值不变。

A1 的领域中位质量分分别为 arxiv 0.782、stackexchange 0.715、commoncrawl 0.707、wikipedia 0.622、github 0.577、book 0.569。c4 的高冲突率为 0.2876；其中“洁净规范高而语言表达低”等方向占主导，驱动指标包括行末标点、专业性和教育价值。该分析指出指标结构的不一致，未将相关性解释为文本语义的因果成因。A2/A3 中 arxiv 和 github 的中位分相对 A1 分别变动 $-0.00532$ 和 $-0.00086$。另从 A1 七域按分位和冲突程度抽取 49 条原文，核对 ID、领域与长度；该样本仅是可追溯的定性审计，不等于独立专家盲评。
\fig{problem1/quality/figures/quality_by_domain.pdf}{0.82}{A1 各领域的相对质量分布。分数仅在当前指标与标尺下可比。}
\fig{problem1/quality/figures/conflict_profile_by_domain.pdf}{0.82}{各领域高冲突比例及结构，阈值由 A1 固定。}

\subsection{17 域配比响应面与模型选择}
配比处于单纯形，直接对各占比做无约束回归会引入和为 1 的共线性。借鉴组成数据分析\cite{aitchison}，将正配比以等距对数比变换 $z=\operatorname{ILR}(p)$ 映射到 16 维，再以 Ridge 稳定估计平均 Loss 响应：
\begin{equation}
\widehat L_{\rm mix}(p)=b_0+b^\top z+z^\top H z,\qquad H=H^\top .
\end{equation}
模型仅在 A4--A5 内部五折交叉验证中选择。二次模型的 MSE 为 0.1471，线性模型为 0.1944；留出的 A6--A7 的 1M 平均 Loss RMSE 为 0.23195，Spearman 为 0.66589。该误差针对每个配方的 13 域平均 Loss；逐域合并 RMSE 为 0.33256，不能混为同一指标。真实 $R^2=1-\mathrm{SSE}/\mathrm{SST}$，独立 1M 测试为 0.3023，不能引用早期误算的 0.946。

A16 的 17 域中只有 6 域可直接或近似对应 A1 的评分域。加入已覆盖域质量特征虽然降低内部 CV 误差，却未同时满足预设的独立测试 2\% 改善门槛，因此最终预测器不加入该质量项。其余 11 域没有独立观测质量分，不以推断代理冒充实测。
\fig{problem1/mixture/figures/mixture_prediction_1m.pdf}{0.83}{1M 独立配比检验：预测与观测的平均 Loss。}

\subsection{推荐配比、非加性与尺度限制}
主推荐取实测标准化 Loss 最优前 10\% 的 52 个配方的质心。占比最高的五域为 pile\_cc 15.13\%、wikipedia\_en 11.15\%、pubmed\_central 10.39\%、stackexchange 8.54\% 和 uspto\_backgrounds 7.99\%。其和为 53.20\%，单域最大份额为 15.13\%。这是下一轮训练的实验候选，不是已验证的全局最优。响应面边界优化解单列为诊断结果，不替代质心。
\fig{problem1/mixture/figures/recommended_mixture.pdf}{0.84}{17 域实测优良配方质心及其占比。}

为讨论组合效应，保存二次 ILR 响应面的系数、Hessian 和按配方行重抽样结果；在训练均值及推荐质心两个锚点，对无序领域对做各增加 0.5 个百分点的联合扰动。136 对无序组合经 BH--FDR 修正后，分别有 30 与 38 对呈模型内非加性。结论依赖锚点、重分配路径和二次响应面，不能称为领域的因果互补。对应的 272 条有向一阶替换仅作局部预测，不赋予独立显著性。

A6--A11 的成对检验表按 1M、60M、1B 分层：60M 与 1B 的平均 Loss 原值 RMSE 分别为 1.5364、2.7544，含有很大的尺度截距偏差；逐尺度去均值后的 RMSE 分别为 0.2160、0.1537，排序相关为 0.6258、0.5266。故 1M 模型移到 60M、1B 后只能部分保持配比效应；在 10B、70B 的估计集上排序相关转负，而这些 Loss 本身又是估计值。因此不用 1M 响应面确定大模型配比，也不将其绝对 Loss 与后续 B6 相加。

\section{问题二：质量修正标度律与跨问接口}
\subsection{参数量--数据量基线}
B1 的 \texttt{pythia\_training\_log\_existing.csv} 有 1176 个检查点，8 条参数规模轨迹各 147 点；B3 的 8 个插值文件共 4000 行且 \texttt{interpolated=1}，只用于连续性审计。以 $N_B,D_B$ 为十亿单位，拟合
\begin{equation}
L_0(N_B,D_B)=E+A N_B^{-\alpha}+B D_B^{-\beta}.
\end{equation}
得到 $E=1.6898,A=0.35398,B=1.24031,\alpha=0.33998,\beta=0.27988$。按完整参数规模轨迹留出，RMSE 为 $0.000151$。这异常小，说明 B1 几乎遵循规则化幂律；轨迹 Bootstrap 给出的极窄区间主要反映该数据集内变化，不能视为真实预训练项目的完整认知不确定性。B4 的 57 条和 B5 的 44 条外部记录因模型族、分词器和验证集不同，原始 Loss 存在协议截距差；去均值 RMSE 分别为 0.1919 和 0.1965，不能仅用 B1 内部误差宣称跨协议精度。B2 的半合成模型族外记录与 B3 的插值轨迹另作压力审计；B9 的 132 条大模型元数据界定远域规模，B10 的 128 条 Loss 为估算值，不计作百亿参数以上真实精度验证。经典标度律及算力最优思想参见\cite{kaplan,hoffmann}。
\fig{problem2/figures/m0_scaling_fit.pdf}{0.86}{B1 训练检查点与参数量--数据量基线拟合。}

\subsection{质量作用位置的竞争模型}
B6 为 360 条半合成质量实验。对质量仅作用于数据项、仅作用于参数项以及同时作用于两项等结构做完整 $(N,D)$ 单元分组验证，并将 $D_B=600$ 与观测上限 300B 内的验证分开。分组 CV 选中的 \texttt{exp\_both} 为
\begin{equation}
L_1(N_B,D_B,Q_B)=\delta_B+
e^{\gamma_Q(1-Q_B)}
\left(A N_B^{-\alpha}+B D_B^{-\beta}\right)+E,
\label{eq:m1}
\end{equation}
其中 $Q_B\in[0.1,1]$，$\delta_B$ 是 B6 来源截距。所选结构的 $\gamma_Q=0.3650$，D$\leq$300B 的分组 CV RMSE 为 0.05951；无质量项为 0.14005。$D_B=600$ 的单列审计 RMSE 为 0.05576，B7 新质量水平且 $D_B\leq300$ 的 RMSE 为 0.04459。它们检验的是半合成体系内的一致性，不能转写为真实训练提质效果。B8 相邻 $Q$ 提升仅 15.64\% 对应 Loss 不升，与 B6/B7 的约 75\% 相反，故不参与主拟合，也不反转其 $Q$ 定义。不同生成机制只是可能解释，尚无原始生成脚本证实。
\fig{problem2/figures/m1_cv_comparison.pdf}{0.82}{质量修正的竞争结构在 B6 分组验证中的误差。}

\subsection{边际效用、参数等价量及 IsoFLOP}
记 $X=A N_B^{-\alpha},Y=B D_B^{-\beta},G=e^{\gamma_Q(1-Q_B)}$，则
\begin{equation}
\frac{\partial L_1}{\partial N_B}=-\frac{\alpha XG}{N_B},\quad
\frac{\partial L_1}{\partial D_B}=-\frac{\beta YG}{D_B},\quad
\frac{\partial L_1}{\partial Q_B}=-\gamma_Q(X+Y)G .
\end{equation}
固定 $D_B,p$，将 $Q_B$ 提高 0.1 与增加 $N_B$ 的模型预测 Loss 改善相等，通过一维求根而非误用只修正数据项的闭式求解。B6 支持网格内，参数倍数 $N'_B/N_B$ 的中位数为 1.2456，第 5 与第 95 百分位分别为 1.1524 和 1.5661；其数值随 $N_B,D_B$ 改变，不是统一换算系数。以上均是选中模型及 B6 标尺下的条件等价。
\fig{problem2/figures/quality_parameter_equivalence.pdf}{0.83}{在固定数据量时，$Q_B$ 提高 0.1 的条件参数等价倍数。}

用 B1 实测算力字段作观测 IsoFLOP 分桶；模型前沿则在经核验的 $C\simeq6ND$ 换算及 B1 支持域内优化。由于式(\ref{eq:m1})的质量因子同时乘两项，固定 $C$ 时模型内的最优 $N,D$ 不随 $Q_B$ 改变，质量只移动预测 Loss 前沿。这是公式的条件性质，不是额外训练实验。
\fig{problem2/figures/isoflop_frontier.pdf}{0.86}{B1 算力支持域内的观测分桶与选中质量模型的条件前沿。}

\subsection{配比项与质量标尺的可识别性}
第一问保存的局部配比响应可作条件接口
\begin{equation}
L_{\rm cand}=L_1(N_B,D_B,Q_B)+w_p(N_B)\Delta_p(p).
\end{equation}
但 A4 配比 Loss 与 B6 质量 Loss 无共同实验协议，故加性、乘性及 $Q_B\times p$ 项的优劣尚不能由附件直接检验；上式仅是候选结构，不用于第三问绝对预测。A16 仅为 6/17 域提供直接或近似质量映射；用代理补齐 17 域后，推荐配方相对训练均值的 $Q_A$ 变化为 $-0.00305$。仅根据已映射域、将未映射域质量限制在 $[0,1]$ 时，变化界为 $[-0.05749,0.03622]$，跨过零；这不是置信区间。没有 $Q_A$ 与 $Q_B$ 的成对实验标定，故不把 $Q_A$ 代入式(\ref{eq:m1})。

\fig{problem2/figures/quality_bridge_identification.pdf}{0.82}{两组配方质量代理差及仅凭已映射域得到的部分识别界。}

问题一的逐域尺度指数仅在具有 Loss 目标的 13 域可估。以 1M 为参考点的 $w_i(N)=(N/10^6)^{-\kappa_i}$ 中，pile\_cc、hackernews、gutenberg\_pg\_19 的 $\kappa_i$ 分别为 $-0.232,-0.149,-0.063$，其权重随规模放大，不能概括为所有配比效应都衰减；超过 1B 的敏感性分析若改以 1B 为参考点，须重新标明 $N_0$。

\section{问题三：给定预算的条件资源配置}
\subsection{决策变量、支持域与成本}
将第一问的 17 域实测优良配方质心固定为外生候选 $p^{\rm ref}$。根据题目允许以前问配比作固定输入的要求，优化变量取 $(N_B,D_B,Q_B)$；因 A4 与 B6 没有联合实验，不能把此计算称为 $N,D,Q,p$ 四变量联合最优。模型响应采用式(\ref{eq:m1})，限定在 $N_B\in[0.070542,11.965825]$、$D_B\in[10,299.893]$、$Q_B\in[0.4,1]$。连续 $Q_B$ 的结果是 B6 离散水平间插值。

总成本按照题给附录 B 写为
\begin{equation}
C_{\rm tot}=6ND+\eta ND L_{\rm ctx}
+D\,[g(Q_B)-g(Q_{B,0})]_+,\qquad
g(Q_B)=10^7e^{6Q_B},
\label{eq:cost}
\end{equation}
主情景取 $\eta=2\times10^{-4}$、$L_{\rm ctx}=2048$、$Q_{B,0}=0.4$。式中 $N,D$ 是原始参数与 token 数；三个成本项均为 FLOPs。$g$ 的指数、幂及对数备选形式均是题设情景而非市场报价。求解 $\min L_1$，约束 $C_{\rm tot}\leq C$ 及前述支持域；逐候选回代成本、边界和目标值，并检查角点的一阶必要条件。

\subsection{外生上下文的临界值}
由式(\ref{eq:cost})的基础训练与注意力两项相等，得到
\begin{equation}
6ND=\eta NDL_{\rm ctx}\quad\Longrightarrow\quad
L_{\rm ctx}^{\rm crit}=\frac{6}{\eta}=30000.
\end{equation}
C7 的最大位置容量记录覆盖 2048、4096、8192、32768 和 131072；这些是模型可支持窗口，并非训练时实际采用的序列长度。32768 情景略高于临界值，此时注意力代理与基础训练之比为 $\eta L_{\rm ctx}/6\approx1.092$。因此本文在 C7 可行容量上做情景敏感性，不把 2048 视为所有模型的实测上下文。
\subsection{主情景结果与边界解释}
\begin{table}[htbp]\centering\caption{指数质量成本及 2048 上下文下的条件配置}
\small\begin{tabular}{rrrrrl}\toprule
预算 FLOPs & $N_B^*$ & $D_B^*$ & $Q_B^*$ & 预测 Loss & 活跃边界\\\midrule
$10^{19}$ & 0.12898 & 10.000 & 0.55744 & 3.3064 & $D$ 下界\\
$10^{21}$ & 2.3231 & 53.152 & 1.0000 & 2.3802 & $Q$ 上界\\
$10^{22}$ & 5.6433 & 249.410 & 1.0000 & 2.1676 & $Q$ 上界\\
$10^{24}$ & 11.966 & 299.893 & 1.0000 & 2.1100 & $N,D,Q$ 上界\\\bottomrule
\end{tabular}\end{table}
\fig{problem3/figures/budget_path.pdf}{0.87}{预算路径与活跃约束；高预算平台由观测支持域截断。}

将结构性转移预先定义为最优解中 $Q_B$ 的下界、内点、上界状态或 $N_B,D_B$ 支持域边界的活跃集合改变。沿预算网格扫描并对活跃集变更作二分定位：指数成本主情景约在 $8.94\times10^{19}$ FLOPs 脱离 $D$ 下界，在 $1.54\times10^{20}$ FLOPs 使质量达到上界；约在 $1.34\times10^{22}$ 和 $2.42\times10^{22}$ FLOPs 先后触及 $D,N$ 上界。这些阈值是当前模型和支持域的数值转折，不是现实训练系统的普适算力门槛；连续成本份额的平滑变化也不称为结构转移。
$10^{19}$ FLOPs 的连续最优质量 0.55744 不属于 B6 的原生离散水平；在其实际水平上重新优化 $N,D$ 得 $Q_B=0.6$，模型预测 Loss 增加 0.00524。此预算的 $D_B=10$ 恰好触及下界，将数据下界向内提高到 20B 时预测 Loss 从 3.3064 升到 3.4199，表明低预算结果受边界设定控制。$10^{21}$ FLOPs 下，用题设幂/对数质量成本替换指数成本，预测最优数据量分别为 48.923B/68.765B；不能把其中一式单独当作真实报价。

在 $10^{24}$ FLOPs 情景，最优点触及三项支持上界，仍有 97.582\% 预算未用。该“最优”只在现有训练数据范围内成立；不能据此认为现实高预算不值得继续扩大模型或数据。成本参数及 B6 质量机制的不确定性没有被窄 Bootstrap 区间完整覆盖。第一问配比质心相对训练均值的 1M 响应面预测差为 $-0.00623$，配方行 Bootstrap 的 95\% 区间为 $[-0.01715,0.00316]$，并跨过零；该差不能加入表中的 B6 绝对 Loss。

\section{问题四：开放模型能力与算力放缓情景}
\subsection{样本定义和三个可比性边界}
采用 C1 六任务同协议均分 $S$，按 C2 中权重明确为 Yes 且许可证属于预设宽松集合来构造操作性“严格开放”样本。原始 4576 条经日期、参数量与六任务筛选及模型去重后，严格开放为 236 个模型，其中预训练 42 个、后训练 141 个；其余类型不混入分层主回归。此筛选可复算，但不等于法律上的开源认证或提交当日已开放证明。
\fig{problem4/figures/q4_sample_attrition.pdf}{0.82}{题目附件 C1/C2 到严格开放样本的筛选过程。}

C1/C4 的算力连接须满足名称、来源身份、参数、发布时间、领域和权重状态的一致性；得到 29 个含算力的预训练连接，其中外部来源可支持的敏感性子集为 15 个、6 个家族。该子集最高算力 $2.94\times10^{23}$ FLOPs，仍低于后文的情景起点。C4 的数据量与权重字段按缺失和口径一致性审计；能与严格开放后训练得分可靠成对连接的训练算力记录为零，故不将 C4 数据量缺失值补成训练 token 数，也不拟合后训练层的联合 $(N,D,C)$ 能力模型。外部来源核查只对连接做补充敏感性，不改写题目附件主字段。C3 的 4573 条排行榜行可匹配 C1；另 26 条早年文献补录的 Average 与六任务均值口径不一致，不能拼为 2019 年以来的同协议长时间序列。C8 的 1954 个 JSON 可做 BBH 叶任务汇总，但汇总分数与 C1 的 BBH 数值标尺未证实相同。以上共同构成样本、跨表和评分的硬边界。

\fig{problem4/figures/q4_observed_frontier.pdf}{0.83}{严格开放模型在 C1 同协议下的观测前沿，按模型类型分层。}

\subsection{规模与时间的条件分解}
将六任务均分映射到 logit 空间，对 $\log_{10}N$ 及排行榜提交时间用 Ridge 拟合；按模型家族分组选择正则化，2025 年 1--3 月作未来时段留出，以家族等权误差评估。严格后训练层中，仅参数量模型家族均衡 RMSE 为 8.0189，加入时间后为 11.736，差值 3.717 的家族 Bootstrap 95\% 区间为 $[-1.113,7.535]$，同时覆盖改善与恶化，故既不能证明时间项更优，也不能在未预设等价界值时宣称两模型等价。严格预训练晚期样本仅 5 个模型，也不足以支撑长期趋势。

在早晚期共同参数支持范围内，对严格后训练层做两因素 Shapley 描述分解：规模项 $-4.876$ 分、条件时间残差 $+5.768$ 分，模型合计 $+0.892$ 分，而样本观测均分变化为 $+2.316$ 分。模型总变化的家族 Bootstrap 区间跨 0；时间项混合架构、数据、后训练和样本选择差异，不是纯技术进步贡献，也不能计算稳定的因果比例。剔除 MATH Lvl 5 后早晚差为 1.173 分，而该单任务差为 14.269 分，显示综合均分对任务构成敏感。
\fig{problem4/figures/q4_conditional_decomposition.pdf}{0.83}{共同参数范围内的描述性规模与时间残差分解。}

\subsection{前沿分位与综合分数的补充检验}
题目关注能力前沿，均分回归不能直接替代条件上分位。为此在严格开放后训练层尝试 $0.9$ 条件分位数回归（它不是最高值），仍以 2025-01 前 74 个模型训练，按 14 个家族分组选择正则化参数，并将 2025-01 至 03 月的 67 个模型（6 家族）整体留出。常数、仅参数量及参数量加时间三类候选的训练期家族均衡 Pinball 损失分别为 2.104、1.803、1.713；训练期选中含时间模型。但在未来时段，其家族均衡损失为 1.322，高于仅参数量的 1.073；逐模型经验覆盖率分别为 0.970 和 0.881，训练分组验证中选中模型的逐模型覆盖率仅 0.689。分位水平为 0.9 不意味着这些异质家族样本具有 90\% 覆盖保证。该补充模型只检验短期前沿口径，不替换正文的描述性分解，也不生成 12/24 个月分数。
\fig{problem4/figures/q4_frontier_quantile_holdout.pdf}{0.86}{条件 90\% 分位候选在未来时段留出上的预测与经验覆盖率；仅有 6 个留出家族。}

六任务均分保留为 C1 的同协议主指标。另以训练期各任务中位数及四分位距固定标准化标尺，早期至晚期的六任务平均变化为 0.202 个训练期四分位距单位；剔除 MATH Lvl 5 后为 0.118。原始单位下六任务均分变化为 3.356 分，二者量纲不同，不直接比较数值大小。MATH Lvl 5 的单任务标准化变化为 0.622，仍提示综合分数对该任务敏感。固定训练期标尺避免用未来样本重设刻度，但没有消除模型家族组成的变化。

在 C1/C4 的 29 个预训练连接上，探索性家族留出比较给出仅算力与参数量加算力模型的 RMSE 分别为 6.691 和 7.117；在来源支持的 15 个连接上分别为 5.178、5.112。结果方向不稳定，且严格开放后训练连接为零，故不能据此分配算力与技术的因果贡献。C6 的中可比 68 点分散在 30 个原始来源标签、其中 14 个标签只有一个点；将它们扩充到高可比 7 点的桥接式会引入未识别的协议偏差，故维持原桥接边界。
\subsection{Loss 桥接及 12/24 个月情景}
C6 与标度律高可比的仅有同一家族 Pythia 的 7 个点。逐点留一的单调 Loss--评分映射 RMSE 为 0.4366，尚差于常数基线 0.4244，且完全没有跨家族验证。因此不能把问题三的 Loss 优化直接换算成排行榜分数。C4 以截至 2025-03-13 的 2025 年不完整年份中位算力 $5.7144\times10^{23}$ FLOPs 为参考，按题给“算力放缓”构造停滞、历史增速减半等条件路径。12 个月的条件算力为 $5.7144\times10^{23}$ 至 $1.0610\times10^{24}$ FLOPs，24 个月为 $5.7144\times10^{23}$ 至 $1.9699\times10^{24}$ FLOPs；历史增速中点仅作对照。改用 2024 年完整年份中位数，所有路径降为上述数值的 0.3323 倍，说明起点选取高度影响情景。

C1/C4 可靠连接的 29 个算力--得分点全部是预训练模型，严格开放后训练的成对点为零。基于这一断裂以及未来留出表现，12/24 个月的具体能力分数在附件下不可识别。可将桥接断点写为 $S_f=h(L)+u_f$，其中 $u_f$ 是家族特有的协议与能力偏移。C6 高可比数据只识别 Pythia 家族内的局部关系，不约束未来新家族的 $u_f$。对任意 $s\in[0,100]$，均可选取一个新家族偏移使 $S_f=s$ 而不改变现有 C6 观测；故在不加跨家族可迁移假设时，新模型分数的识别集为 $[0,100]$。若前沿定义为已有纪录与未来分数之最大值，识别集才缩为 $[46.889,100]$，并非统计置信区间。
图中仅展示算力假设和识别状态，不绘制虚构能力曲线。
\fig{problem4/figures/q4_compute_scenario_identification.pdf}{0.85}{12/24 个月条件算力路径及能力分数不可识别的审计结果。}

\section{模型检验、敏感性与适用范围}
\subsection{分层验证}
问题一的配比模型以 A4--A5 训练、A6--A7 的 1M 记录作独立检验；A12--A15 大尺度估计仅用于压力测试。问题二 M0 按整条参数规模轨迹留出，M1 按完整 $(N,D)$ 单元分组，$D_B=600$ 和 B7 新质量水平单列，避免同一单元同时进入拟合与验证。问题三对各预算回代原始 FLOPs、支持域与活跃约束，比较三种题设质量成本、基线质量、上下文长度及 B6 离散质量档位。问题四按家族分组和时间留出，另核查开放状态、C1/C4 连接、历史评分协议及 C8 叶任务标尺。

\begin{table}[htbp]\centering\caption{关键验证数值与解释边界}
\small\begin{tabular}{p{3.2cm}p{2.8cm}p{5.7cm}}\toprule
验证 & 结果 & 解释\\\midrule
配比 1M 独立测试 & RMSE 0.23195 & 13 域平均 Loss，不能代表大尺度绝对预测\\
B1 留完整轨迹 & RMSE 0.000151 & 附件内部强规则关系，不代表现实泛化误差\\
B6 单元分组 CV & RMSE 0.05951 & 半合成质量响应的结构选择\\
B7 新质量档位 & RMSE 0.04459 & 限于 B7 的 D$\leq$300B\\
后训练未来留出 & 8.0189 / 11.736 & 仅规模 / 加时间的家族均衡 RMSE\\\bottomrule
\end{tabular}\end{table}

\subsection{共同限制与不得外推之处}
质量评分 $Q_A$ 是按 A1 标尺构造的相对指标，$Q_B$ 是 B6 实验水平；二者没有共同样本、共同 Loss 或共同配比的联合标定。第三问固定的 $p^{\rm ref}$ 因此不改变数值目标和成本，表中仅为 $(N,D,Q_B)$ 的条件最优。问题一二次 ILR 的非加性检验提供模型内局部证据，不证明领域内在互补。问题二逐域规模指数有负值，且只有 13 域目标可估，跨尺度统一衰减不成立。问题四排行榜记录、算力元数据、Loss 测试之间存在样本选择和协议差异，无法构建经验证的长期能力前沿。

为避免不确定性过窄，文中将三类量分开：数据重抽样区间描述当前附件内的不稳定性；支持域扫描描述边界选择；成本函数与质量机制替换描述结构情景。三者不能相加为单一“95\% 预测区间”。尤其 B1 的极窄轨迹 Bootstrap、第三问的条件 Bootstrap 与第四问的宽逻辑界，性质完全不同。

\section{结论与下一步验证}
本文给出一条以题目附件为主、在证据断点处停止数值传递的资源配置分析链。第一问的稳健评分揭示领域间相对差异和 c4 指标冲突；二次 ILR 对 1M 配比 Loss 有可检验预测力，但实测优良配方质心仍是待训练验证的候选。第二问的 B1 幂律和 B6 质量修正分别在各自数据内成立；选中结构表明半合成质量水平可改变参数与数据项的条件 Loss，但不能证明第一问质量评分具有相同效应。第三问在固定候选配比和题设成本下得到预算路径；低预算受 $D$ 下界影响，高预算出现支持域饱和。第四问可描述严格开放样本的观测能力和条件算力路径，却不能可靠估计 12/24 个月能力数值。

最直接的后续实验是：在同一训练与评测协议下，对多个 $N,D$、至少两种 17 域配比、多个可复测的质量水平作成对训练；同时记录真实质量处理成本与上下文长度。这样才能标定 $Q_A\leftrightarrow Q_B$、比较配比与质量的交互结构，并把前三问的 Loss 决策桥接至第四问能力指标。没有这些共同实验前，本文不把条件最优和历史关联包装为现实的因果或长期预测结论。


\begin{thebibliography}{9}
\bibitem{problem} 第二十三届中国研究生数学建模竞赛 F 题. 算力约束下提升大语言模型能力的资源配置建模及随题附件.
\bibitem{aitchison} Aitchison J. The Statistical Analysis of Compositional Data. Chapman and Hall, 1986.
\bibitem{kaplan} Kaplan J, et al. Scaling Laws for Neural Language Models. arXiv:2001.08361, 2020.
\bibitem{hoffmann} Hoffmann J, et al. Training Compute-Optimal Large Language Models. NeurIPS, 2022.
\end{thebibliography}
\appendix
\section{复算索引与论文使用说明}
论文正文数值由仓库中的四份结果报告及其 CSV/JSON 表核对；程序分别为 \texttt{code/problem1/problem1.py}、\texttt{code/problem2/problem2.py}、\texttt{code/problem3/problem3.py}、\texttt{code/problem4/problem4.py}，输出分别在 \texttt{outputs/problem1} 至 \texttt{outputs/problem4}。每问的独立验收程序及报告位于对应的 \texttt{code/problemN/verify\_problemN.py} 和 \texttt{reports/问题N/验证验收/}。论文引入的 PDF 图均引用这些程序已生成的文件；论文未重新构造实验点或补填缺失评分。第四问另有前沿分位与任务标尺补充程序 \texttt{code/problem4/problem4\_frontier\_sensitivity.py}，其独立回代脚本为 \texttt{code/problem4/verify\_problem4\_frontier\_sensitivity.py}，审查意见核查见 \texttt{reports/问题四/验证验收/IMAGE\_SUGGESTION\_AUDIT.md}。

正文保留三类待核验信息：参赛队号、学校及队员信息不在数据附件中，因此当前匿名稿不填虚构署名；正式提交前按竞赛当届规则处理封面。题目附件中的版权、来源和模型开放许可仍应按最终提交要求再次人工核查。若重跑数据与现有报告不一致，应先更新结果报告，再同步修改论文表格与图注，而非只改文字。

\section{主要符号}
\begin{tabular}{p{2.1cm}p{10.5cm}}\toprule
符号 & 含义\\\midrule
$Q_A$ & A1 指标体系下的相对质量评分\\
$Q_B$ & B6 原生质量水平，本文只在其观测支持内使用\\
$p$ & 17 域占比向量，$\sum_i p_i=1$\\
$N_B,D_B$ & 以十亿为单位的参数量和训练 token 数\\
$L_0,L_1$ & B1 基线标度式及 B6 来源校准后的质量式\\
$C,L_{\rm ctx}$ & FLOPs 预算及外生上下文长度\\
$S$ & C1 同协议六任务平均能力分数\\\bottomrule
\end{tabular}

\section{AI 工具使用披露}
本项目使用 OpenAI Codex 协助梳理题目与数据字段、讨论模型方案、编写和修改分析程序、运行结果核验、绘制图表以及起草本论文文字。论文中的数值和图形须能由题目附件及仓库程序复算；Codex 输出本身不作为数据来源或学术依据。参赛队在正式提交前须逐式、逐表和逐图独立核验，并对模型选择、计算过程与最终论证负责。

\end{document}
